\section{Products of resolvents with independent complex powers}
\label{mr:section}

Question 10 of the incoming \emph{Quartic Tails and Complex Resolvents}
report asks for products of several nonreal cotangent resolvents, including
parameters in opposite half-planes \cite{repo:quartic}. We answer its
analytic and transformation parts for any finite number of factors.
The answer has two distinct features: every Fourier coefficient has a
geometrically convergent Gamma-series representation, and the complete
endpoint subtraction depends on a signed sum of exponents as well as on
their total. The latter dependence produces genuine extra poles when both
half-planes occur.

\subsection{Branches and notation}

Let $z_1,\ldots,z_d\in\C\setminus\R$ and
$\boldsymbol\lambda=(\lambda_1,\ldots,\lambda_d)\in\C^d$. Repeated
values of the $z_j$ are allowed. Define
\begin{align}
 K(x)&=\pi\cot(\pi x),& R_z(x)&=(K(x)-z)^{-1},\notag\\
 \epsilon_j&=\sgn\Im z_j,&
 q_j&=\frac{z_j-\epsilon_j\pi\ii}{z_j+\epsilon_j\pi\ii},&
 C_j&=-\frac1{z_j+\epsilon_j\pi\ii},\label{mr:diskdata}\\
 a&=\sum_{\epsilon_j=1}\lambda_j,&
 b&=\sum_{\epsilon_j=-1}\lambda_j,&
 \Lambda&=a+b,\qquad E=a-b.\notag
\end{align}
In particular $|q_j|<1$. In every occurrence of $C_j^{\lambda_j}$ use
\[
 \Log C_j=\Log(1-q_j)-\log(2\pi)+\epsilon_j\pi\ii/2,
\]
where $\Log(1-q_j)$ is holomorphic on the disk and vanishes at $q_j=0$.
The continuous logarithm of $R_z$ on $(0,1)$
is fixed by $\Log R_z(x)-\log x\to0$ as $x\downarrow0$. For
$\epsilon=\sgn\Im z$, its imaginary part lies between $0$ and
$\epsilon\pi$. Powers in
\begin{equation}\label{mr:kernel}
 P_{\boldsymbol\lambda}(x)
       =\prod_{j=1}^d R_{z_j}(x)^{\lambda_j}
\end{equation}
use these logarithms. Put
\[
 B(\boldsymbol\lambda,\mathbf q)
       =\prod_j(1-q_j)^{\lambda_j},
 \qquad
 V(x;\boldsymbol\lambda)
       =\prod_j\bigl(\pi x\cot\pi x-z_jx\bigr)^{-\lambda_j}.
\]
The powers defining $B$ use the holomorphic logarithms on $|q_j|<1$.
The powers defining $V$ use their analytic logarithms at $x=0$, where
$V(0)=1$. Direct algebra gives both the global expression
\begin{align}
 P_{\boldsymbol\lambda}(x)
  ={}&\pi^{-\Lambda}B\,\sin^{\Lambda}(\pi x)e^{\pi\ii E x}
       \prod_{\epsilon_j=1}(1-q_je^{2\pi\ii x})^{-\lambda_j}
       \prod_{\epsilon_j=-1}(1-q_je^{-2\pi\ii x})^{-\lambda_j},
       \label{mr:sinefactor}
\end{align}
and the two endpoint expansions
\begin{equation}\label{mr:endpoints}
 P_{\boldsymbol\lambda}(x)=x^\Lambda V(x;\boldsymbol\lambda),\qquad
 P_{\boldsymbol\lambda}(1-y)
       =e^{\pi\ii E}y^\Lambda V(-y;\boldsymbol\lambda).
\end{equation}
The leading coefficient at zero is exactly one. This normalization
will fix every contact term below.

\subsection{An explicit Gamma series for the full Fourier spectrum}

Define the exponentially decreasing coefficient sequences
\begin{equation}\label{mr:UV}
 U_n=[w^n]\prod_{\epsilon_j=1}(1-q_jw)^{-\lambda_j},\qquad
 W_n=[w^n]\prod_{\epsilon_j=-1}(1-q_jw)^{-\lambda_j}.
\end{equation}
An empty product has coefficient $1$ at $n=0$ and $0$ otherwise.
Every coefficient is a finite sum; for example
\[
 U_n=\sum_{\substack{(n_j)\ge0\\\sum_{\epsilon_j=1}n_j=n}}
          \prod_{\epsilon_j=1}\frac{(\lambda_j)_{n_j}q_j^{n_j}}{n_j!}.
\]
For $k\in\Z$ set
\begin{equation}\label{mr:Gk}
 G_k(a,b)=
 \frac{e^{\pi\ii(a-b)/2}(-1)^k\Gamma(1+a+b)}
 {(2\pi)^{a+b}\Gamma(1+a-k)\Gamma(1+b+k)}.
\end{equation}
Reciprocal Gamma factors are read as entire functions. Restrictions
through an intersection of a numerator pole and a reciprocal-Gamma zero
are made as meromorphic restrictions, never as an informal $0\cdot\infty$.

\begin{theorem}[Fourier coefficients of any finite product]
\label{mr:fourier}
For $\Re\Lambda>-1$, the ordinary Fourier coefficients of
\eqref{mr:kernel} are
\begin{equation}\label{mr:ck}
 c_k(\boldsymbol\lambda,\mathbf z)
   :=\int_0^1P_{\boldsymbol\lambda}(x)e^{-2\pi\ii kx}\dd x
   =B\sum_{n,m\ge0}U_nW_mG_{k-n+m}(a,b).
\end{equation}
The double series converges normally, geometrically in the two summation
indices on compact subsets of $|q_j|<1$, away from its possible parameter
poles. It supplies a joint meromorphic continuation with possible poles
only on $\Lambda=-1,-2,\ldots$. Its mean is $c_0$.

If every $z_j$ is in the same half-plane, then
\begin{equation}\label{mr:samemean}
                         c_0=\prod_jC_j^{\lambda_j}.
\end{equation}
For mixed half-planes, \eqref{mr:ck} gives an explicit two-sided
spectrum and \eqref{mr:samemean} generally fails.
\end{theorem}

\begin{proof}
The classical Fourier beta integral \cite[Eq.~5.12.6]{dlmf:beta} reads
\[
 \int_0^1\sin^\Lambda(\pi x)e^{\pi\ii\nu x}\dd x
  =\frac{e^{\pi\ii\nu/2}\Gamma(1+\Lambda)}
         {2^\Lambda\Gamma(1+(\Lambda+\nu)/2)
                         \Gamma(1+(\Lambda-\nu)/2)},\qquad
 \Re\Lambda>-1.
\]
Substitution $\nu=a-b-2k$ gives $G_k$. Expanding only the analytic
factors in \eqref{mr:sinefactor}, whose series converge uniformly on
the circle, and integrating against the integrable endpoint factor
proves \eqref{mr:ck}.

Choose a radius $\rho>1$ with $\rho\max_j|q_j|<1$. Cauchy's
estimate bounds $U_n,W_n$ by compact-set constants times $\rho^{-n}$
and $\rho^{-m}$. Gamma reflection and the Gamma-ratio asymptotic show
that $G_k$ has at most polynomial growth in the integer $k$, uniformly
on compact parameter sets avoiding $\Lambda=-1,-2,\ldots$.
Consequently the series and all fixed parameter derivatives converge
normally there. Formula \eqref{mr:Gk} proves the stated meromorphy.

In a common half-plane put $w=e^{2\epsilon\pi\ii x}$. The kernel is
\[
 \left(\prod_jC_j^{\lambda_j}\right)
       (1-w)^\Lambda\prod_j(1-q_jw)^{-\lambda_j}.
\]
Its holomorphic disk function has constant coefficient $\prod C_j^{\lambda_j}$.
Abel limits are integrable when $\Re\Lambda>-1$, by the combined
endpoint exponent, and give \eqref{mr:samemean}. This argument uses
the total exponent, not separate inequalities on individual exponents.
\end{proof}

\begin{remark}[A finite hypergeometric specialization]
With one factor at $z_+=\pi\ii$ of exponent $a$ and one lower-half-plane
factor of exponent $b$ and disk coordinate $r$, \eqref{mr:ck} gives
\begin{equation}\label{mr:gaussmean}
 c_0=(1-r)^bG_0(a,b)\,
       {}_2F_1\left(\begin{matrix}-a,b\\1+b\end{matrix};r\right).
\end{equation}
This follows from $G_m/G_0=(-a)_m/(1+b)_m$ in the initial
nonexceptional domain. At exceptional parameters the full expression is
continued as a unit. Thus even one opposite orientation changes the
constant-mode formula from a product of constants to a hypergeometric
function. Formula \eqref{mr:ck}, rather than a name for a new function,
specifies the general answer by elementary finite coefficients and Gamma
values.
\end{remark}

\subsection{Hurwitz, Stieltjes, and log-Gamma transforms}

Let $p\ge0$ be an integer and put
\[
 A_p(s)=\prod_{j=1}^p(s-j),\quad A_0(s)=1,\qquad
 \sigma_+=-2\pi\ii,\quad \sigma_-=2\pi\ii,
 \quad \Log\sigma_\pm=\log(2\pi)\mp\pi\ii/2.
\]
Define, initially in an absolute-convergence region,
\begin{equation}\label{mr:Dirichlet}
 D_+(u)=\sum_{n\ge1}\frac{c_n}{n^u},\qquad
 D_-(u)=\sum_{n\ge1}\frac{c_{-n}}{n^u}.
\end{equation}
The definitions retain all coefficients, including their normalization $B$.

\begin{theorem}[The joint spectral identity]
\label{mr:hurwitz}
In the ordinary-integral domain
\[
 \Re\Lambda>-1,\qquad \Re(s+\Lambda)>p,
\]
away from the Hurwitz pole for $p=0$, analytic continuation from an
absolutely convergent Fourier region gives
\begin{align}
 J_p(s,\boldsymbol\lambda)
 &:=\int_0^1\partial_x^p\zeta(1-s,x)P_{\boldsymbol\lambda}(x)\dd x\notag\\
 &=\Gamma(s)\left\{\sigma_+^{p-s}D_+(s-p)
                      +\sigma_-^{p-s}D_-(s-p)\right\}.
       \label{mr:Hurwitz}
\end{align}
The completed Stieltjes generator is
\begin{equation}\label{mr:T}
 T_p(s,\boldsymbol\lambda)=J_p(s,\boldsymbol\lambda)
                                  +\mathbf1_{p=0}\frac{c_0}{s}.
\end{equation}
It is jointly meromorphic. Its only possible polar hypersurfaces are
\begin{equation}\label{mr:poles}
       s+\Lambda=p-j\quad(j\ge0),\qquad
       \Lambda=-1-j\quad(j\ge0).
\end{equation}
For $\Re\Lambda>p$,
\begin{equation}\label{mr:ordinaryStieltjes}
 \int_0^1\gamma_m^{(p)}(x)P_{\boldsymbol\lambda}(x)\dd x
                  =m![s^m]T_p(s,\boldsymbol\lambda).
\end{equation}
\end{theorem}

\begin{proof}
In an initial region $\Re s>p+1$, $\Re\Lambda>-1$, Fourier pairing
with the classical Hurwitz expansion \cite[Eq.~25.11.9]{dlmf:hurwitz}
gives zero for the constant mode and
\[
 \int_0^1\partial_x^p\zeta(1-s,x)e^{\pm2\pi\ii nx}\dd x
       =\Gamma(s)\sigma_\pm^{p-s}n^{p-s}.
\]
Here the differentiated Hurwitz series converges absolutely and
uniformly, and $P_{\boldsymbol\lambda}$ belongs to $L^1(0,1)$.
Thus termwise integration is justified directly by its uniform bound;
no regularity assertion for the endpoint values of $P$ is required.
This proves \eqref{mr:Hurwitz} in that region.

Use the exact holomorphic splitting
\begin{equation}\label{mr:split}
 \partial_x^p\left(\zeta(1-s,x)+\frac1s\right)
      =A_p(s)x^{s-p-1}+h_p(s,x),\qquad
 h_p(s,x)=\partial_x^p\left(\zeta(1-s,1+x)+\frac1s\right).
\end{equation}
It proves normal integrability on the stated ordinary domain and gives
analytic continuation of the Fourier identity to that domain. At the
two endpoints, the monomial exponents are respectively
$s+\Lambda-p-1+j$ and $\Lambda+j$. Subtracting any required finite
number of such monomials leaves a normally convergent integral.
Their integration denominators give exactly \eqref{mr:poles}.
The pole of the uncompleted Hurwitz factor is $-c_0/s$ for $p=0$;
\eqref{mr:T} removes it before Stieltjes coefficient extraction.
The convention
$\zeta(1-s,x)+s^{-1}=\sum_{m\ge0}\gamma_m(x)s^m/m!$
then gives \eqref{mr:ordinaryStieltjes}.
\end{proof}

The apparent coordinate poles of separate Dirichlet or Gamma factors
in \eqref{mr:Hurwitz} may cancel. Equation \eqref{mr:poles} concerns
the complete function, as proved by its endpoint integrals.

\begin{corollary}[An ordinary log-Gamma identity]
\label{mr:loggamma}
For $\Re\Lambda>-1$,
\begin{align}
 \int_0^1\log\Gamma(x)P_{\boldsymbol\lambda}(x)\dd x
 ={}&\frac{\log(2\pi)}2c_0\notag\\
 &+\sum_{\epsilon\in\{+,-\}}
     \frac{(\gamma+\Log\sigma_\epsilon)D_\epsilon(1)
                                     -D_\epsilon'(1)}{\sigma_\epsilon}.
       \label{mr:logGamma}
\end{align}
The prime differentiates the Dirichlet order. This is an ordinary
integral; its logarithmic endpoint is integrable under the stated
condition. The formula follows by differentiating \eqref{mr:Hurwitz}
at $s=1$ and using
$\zeta'(0,x)=\log\Gamma(x)-\tfrac12\log(2\pi)$.
\end{corollary}

For a concrete ordinary evaluation, take the special points
$z_\pm=\pm\pi\ii$ with exponents $2,1$. Their Fourier coefficients are
$\ii(-1,3,-3,1)/(8\pi^3)$ at modes $-1,0,1,2$, respectively.
Consequently \eqref{mr:logGamma} reduces to
\begin{align}
 &\int_0^1\log\Gamma(x)\,
           \frac{\sin^3(\pi x)}{\pi^3}e^{\pi\ii x}\dd x\notag\\
 &\qquad=
 \frac{3\{\gamma+\log(2\pi)\}-\log2}{32\pi^4}
       +\frac{\ii\{12\log(2\pi)-7\}}{64\pi^3}.
       \label{mr:loggammaexample}
\end{align}

\subsection{The full completion at each total-exponent resonance}

The preceding polar arrangement is completed explicitly as follows.
Set
\begin{equation}\label{mr:vandb}
 v_j(\boldsymbol\lambda)=[x^j]V(x;\boldsymbol\lambda),\qquad
 b_j(s,\boldsymbol\lambda)=[x^j]h_p(s,x)V(x;\boldsymbol\lambda).
\end{equation}
These are finite, effectively computable expressions. For instance
\[
 v_0=1,\qquad v_1=\sum_j\lambda_jz_j,\qquad
 v_2=\frac12\left(\sum_j\lambda_jz_j\right)^2
       +\frac12\sum_j\lambda_jz_j^2+\frac{\pi^2}3\Lambda.
\]
For $p+j>0$, the coefficient of $x^j$ in $h_p$ is
\[
 \frac{(-1)^{p+j}(1-s)_{p+j}}{j!}\zeta(1-s+p+j),
\]
and at $p=j=0$ it is $\zeta(1-s)+1/s$. Thus all required local
derivatives use only ordinary zeta derivatives and Stieltjes constants.

Fix $\boldsymbol\lambda^0$ with
$L=\sum_j\lambda_j^0\in\Z$, write
$\boldsymbol\lambda=\boldsymbol\lambda^0+\boldsymbol\eta$ and
$\delta=\sum_j\eta_j$. Define
\begin{align}
 H_{p,\boldsymbol\lambda^0}(s,\boldsymbol\eta)
 ={}&T_p(s,\boldsymbol\lambda^0+\boldsymbol\eta)
 -\mathbf1_{L\le p}\frac{A_p(s)v_{p-L}(\boldsymbol\lambda^0+\boldsymbol\eta)}
                                  {s+\delta}\notag\\
 &-\mathbf1_{L\le-1}
 \frac{\left\{1+(-1)^{-L-1}e^{\pi\ii E(\boldsymbol\lambda^0+\boldsymbol\eta)}\right\}
       b_{-L-1}(s,\boldsymbol\lambda^0+\boldsymbol\eta)}{\delta}.
       \label{mr:H}
\end{align}
False-indicator terms are omitted, not evaluated.

\begin{theorem}[All Stieltjes, argument, and exponent derivatives]
\label{mr:completion}
The function \eqref{mr:H} is holomorphic near $(s,\boldsymbol\eta)=(0,0)$.
For $m\ge0$ and a multi-index $\mathbf k\ge0$,
\begin{align}
 &\FP_x\int_0^1\gamma_m^{(p)}(x)
       \prod_j R_{z_j}(x)^{\lambda_j^0}
                       (\Log R_{z_j}(x))^{k_j}\dd x\notag\\
 &\hspace{20mm}=m!\left(\prod_j k_j!\right)
                    [s^m\boldsymbol\eta^{\mathbf k}]
                             H_{p,\boldsymbol\lambda^0}(s,\boldsymbol\eta).
       \label{mr:alljets}
\end{align}
The finite part uses the original endpoint coordinates $x$ and $1-x$.
The individual exponents $\lambda_j^0$ need not be integers; only
their total is integral in this statement. If the total is not integral,
the corresponding local germ of $T_p$ itself gives all the finite parts.
\end{theorem}

\begin{proof}
At zero, \eqref{mr:split} and \eqref{mr:endpoints} give the exact
families
\[
 A_p(s)v_j(\boldsymbol\lambda)x^{s+\Lambda-p-1+j},\qquad
 b_j(s,\boldsymbol\lambda)x^{\Lambda+j}.
\]
At the other endpoint the smooth Hurwitz factor is $h_p(s,-y)$.
Its product with the kernel has coefficient
$e^{\pi\ii E}(-1)^jb_j(s,\boldsymbol\lambda)$ at $y^{\Lambda+j}$.
The only exponents whose integrated denominators vanish at the base
point are $j=p-L$ and $j=-L-1$. Their complete amplitudes give the
two subtraction terms in \eqref{mr:H}.

For rigor, choose a fixed small endpoint interval $[0,\rho]$, and
subtract sufficiently many of the displayed monomials that the remainders
and their fixed derivatives are integrable uniformly on a parameter
polydisk. Every nonresonant integrated monomial is holomorphic there.
For each resonant exponent $\alpha-1$ use
\[
 \frac{\rho^\alpha}{\alpha}
             =\frac1\alpha+\frac{\rho^\alpha-1}{\alpha}.
\]
The second term is entire; its Taylor coefficients are exactly the
constant terms of the logarithmic cutoff integrals. Subtracting the
first term with its full parameter-dependent amplitude therefore makes
the remainder holomorphic and proves \eqref{mr:alljets}. Contributions
depending on the artificial splitting point cancel between the local
and remaining integrals. The same argument proves the nonintegral case,
where no denominator is resonant at the base point.
\end{proof}

\begin{corollary}[Why opposite half-planes produce an extra pole]
\label{mr:phasepole}
On the hyperplane $\Lambda=-M$, $M\ge1$, the regular-endpoint
residue in \eqref{mr:H} is
\begin{equation}\label{mr:mixedresidue}
       (1-e^{-2\pi\ii b})\,b_{M-1}(s,\boldsymbol\lambda).
\end{equation}
If all factors lie in the same half-plane this vanishes identically on
the hyperplane, so its quotient is removable. With both orientations
present it is generically nonzero. Its vanishing at an isolated integral
parameter vector does not make that multivariable quotient removable.
\end{corollary}
\begin{proof}
On $\Lambda=-M$, $E=\Lambda-2b=-M-2b$, and hence
$1+(-1)^{M-1}e^{\pi\ii E}=1-e^{-2\pi\ii b}$.
In a single orientation $b=0$ or $b=\Lambda$. With two orientations,
$b$ varies independently along the hyperplane. The local coefficient
$b_{M-1}$ may introduce further cancellations; generic nonvanishing
already occurs for $M=1$ and the constant test function.
\end{proof}

\subsection{An exactly computable mixed example}

For $z_+=\pi\ii$, $z_-=-\pi\ii$, write their exponents as
$\alpha,\beta$. The kernel and ordinary mean simplify to
\begin{align}
 P_{\alpha,\beta}(x)
   &=\left(\frac{\sin\pi x}{\pi}\right)^{\alpha+\beta}
                           e^{\pi\ii(\alpha-\beta)x},\label{mr:specialkernel}\\
 M(\alpha,\beta)
   &=\frac{e^{\pi\ii(\alpha-\beta)/2}\Gamma(1+\alpha+\beta)}
      {(2\pi)^{\alpha+\beta}\Gamma(1+\alpha)\Gamma(1+\beta)}.
       \label{mr:specialmean}
\end{align}
For any $c,d\in\C$ with $c+d\ne0$, continuation along
$\alpha=-1+ct$, $\beta=dt$ yields
\begin{equation}\label{mr:pathlimit}
           \lim_{t\to0}M(-1+ct,dt)=-2\pi\ii\frac{c}{c+d}.
\end{equation}
The fixed-coordinate finite part of the specialized integrand is instead
\begin{equation}\label{mr:fixedmean}
       \FP_x\int_0^1(K(x)-\pi\ii)\dd x=-\pi\ii.
\end{equation}
The exact correction is the full regular-endpoint quotient
\[
 \frac{1+e^{\pi\ii(\alpha-\beta)}}{1+\alpha+\beta}.
\]
Its limit on the same ray is $-\pi\ii(c-d)/(c+d)$; subtracting it
from \eqref{mr:pathlimit} gives \eqref{mr:fixedmean} for every ray.
Thus the rays $(c,d)=(1,0),(1,1),(0,1)$ yield continued limits
$-2\pi\ii,-\pi\ii,0$, respectively, for the same specialized
function. This is an exact counterexample to an uncompleted independent
exponent limit, not an error in the earlier single-resolvent theorem.

At a general negative-total hyperplane $\alpha+\beta=-M$, the
coefficient entering the mean residue is explicitly
\begin{equation}\label{mr:residuepoly}
 [x^{M-1}]V(x;\alpha,-M-\alpha)
          =(2\pi\ii)^{M-1}\binom{\alpha+M-1}{M-1}.
\end{equation}
Indeed this coefficient is the residue at zero of
$e^{\pi\ii(2\alpha+M)x}(\pi/\sin\pi x)^M\dd x$; putting
$w=e^{2\pi\ii x}$ converts it to
$(2\pi\ii)^{M-1}\Res_{w=1}w^{\alpha+M-1}(w-1)^{-M}\dd w$.
Consequently the residue of \eqref{mr:specialmean} is
\[
 (1-e^{2\pi\ii\alpha})(2\pi\ii)^{M-1}
                              \binom{\alpha+M-1}{M-1}.
\]
For example, at $\alpha=\beta=-1/2$ it is $2$, and direct
integration of $\pi\csc\pi x$ gives
\begin{equation}\label{mr:cscfp}
               \FP_x\int_0^1\pi\csc(\pi x)\dd x
                             =2\log(2/\pi).
\end{equation}

The dependence on the path can survive even when every transverse
linear path gives the expected zero. At $(-1,-1)$ write
\[
 \alpha=-1+(T+D)/2,\qquad \beta=-1+(T-D)/2.
\]
On the curved path $T=cD^2$ with $c\ne0$, the Gamma quotient
\eqref{mr:specialmean} has limit $\pi^2/c$. Indeed its leading terms
are $\pi^2D^2/T-\pi^2T$. The specialized kernel is
$\pi^2\csc^2\pi x$, whose coordinate finite part is zero.
Its full contact is
$(1-e^{\pi\ii D})\pi\ii D/T$, with the same curved-path limit.
The completed value is again independent of the path.

\begin{corollary}[A finite-part logarithmic family in ordinary constants]
\label{mr:negativeLog}
Let $R_\pm=R_{\pm\pi\ii}$ and put
\begin{align}
 Q(a,b)&=e^{-(a+b)\log(2\pi)+\pi\ii(a-b)/2}
              \frac{\Gamma(1+a+b)}{\Gamma(1+a)\Gamma(1+b)},\notag\\
 H(a,b)&=\frac{-2\pi\ii aQ(a,b)-1+e^{\pi\ii(a-b)}}{a+b}.
       \label{mr:negativeLoggerm}
\end{align}
This quotient is holomorphic at $(0,0)$. For integers $r,s\ge0$,
\begin{equation}\label{mr:negativeLogmoment}
 \FP_x\int_0^1(K-\pi\ii)(\Log R_+)^r(\Log R_-)^s\dd x
                      =r!s![a^rb^s]H(a,b).
\end{equation}
Every such value is a finite polynomial in $\pi\ii$, $\log(2\pi)$,
and ordinary zeta values. In particular,
\begin{align}
 \FP_x\int_0^1(K-\pi\ii)\Log R_+\dd x
       &=\frac{\pi^2}{2}+2\pi\ii\log(2\pi),\notag\\
 \FP_x\int_0^1(K-\pi\ii)(\Log R_-)^s\dd x
       &=\frac{(-\pi\ii)^{s+1}}{s+1}.
       \label{mr:negativeLogexamples}
\end{align}
\end{corollary}

\begin{proof}
At $\alpha=-1+a,\beta=b$, \eqref{mr:specialmean} equals
$-2\pi\ii aQ(a,b)/(a+b)$. Subtracting the complete paired endpoint
term gives \eqref{mr:negativeLoggerm}.
On $b=-a$, reflection gives
$Q(a,-a)=e^{\pi\ii a}\sin(\pi a)/(\pi a)$, so the numerator vanishes
identically. Analytic division and the endpoint proof above justify
all coefficients. The Gamma expansion proves the stated constant class.
For an independent check of the second identity,
$(\Log R_-)'=K-\pi\ii$. The endpoint logarithms are
$\log\eta+o(1)$ and $\log\eta-\pi\ii+o(1)$; the constant term of
their $(s+1)$st-power difference is $(-\pi\ii)^{s+1}$.
\end{proof}

\subsection{All polynomial moments and parameter derivatives}

The same beta formula gives an ordinary exponential moment at the
special points:
\begin{equation}\label{mr:expmean}
 \int_0^1e^{tx}P_{\alpha,\beta}(x)\dd x
 =\frac{e^{\pi\ii(\alpha-\beta)/2+t/2}\Gamma(1+\alpha+\beta)}
 {(2\pi)^{\alpha+\beta}
   \Gamma(1+\alpha+t/(2\pi\ii))
   \Gamma(1+\beta-t/(2\pi\ii))}.
\end{equation}
It is valid for $\Re(\alpha+\beta)>-1$ and every complex $t$.
Away from zeros of its value at $t=0$, its logarithmic derivatives are
\begin{align}
 \kappa_1&=\frac12+\frac{\psi(1+\beta)-\psi(1+\alpha)}{2\pi\ii},\notag\\
 \kappa_r&=-\frac{\psi^{(r-1)}(1+\alpha)
                    +(-1)^r\psi^{(r-1)}(1+\beta)}{(2\pi\ii)^r},
                    \qquad r\ge2.\label{mr:cumulants}
\end{align}
Thus its $n$th moment is $M(\alpha,\beta)Y_n(\kappa_1,\ldots,\kappa_n)$,
where $Y_n$ is the complete exponential Bell polynomial. At Gamma zeros,
differentiate the full reciprocal-Gamma formula \eqref{mr:expmean};
the factored Bell expression is only its meromorphic rewriting.
All mixed exponent derivatives follow by ordinary parameter
differentiation in this domain, and by \eqref{mr:H} at resonances.

At nonnegative integer $\alpha=m,\beta=n$, these formulas give finite
generalized-harmonic expressions at every moment order. Explicitly,
\[
 M(m,n)=
 \frac{e^{\pi\ii(m-n)/2}}{(2\pi)^{m+n}}\binom{m+n}{m},
 \qquad
 \kappa_1=\frac12+\frac{H_n-H_m}{2\pi\ii},
\]
and for $r\ge2$,
\[
 \kappa_r=
 \frac{(-1)^{r+1}(r-1)!}{(2\pi\ii)^r}
 \left\{\zeta(r)-H_m^{(r)}
                 +(-1)^r\bigl(\zeta(r)-H_n^{(r)}\bigr)\right\}.
\]
Substitution into the Bell polynomial is a finite identity involving
only harmonic numbers, zeta values, and the stated prefactor.

Finally, for any fixed-coordinate moment $\mathcal M$ in
\eqref{mr:alljets}, differentiation in a resolvent parameter gives
\[
 \partial_{z_j}\mathcal M(\boldsymbol\lambda,\mathbf k)
 =\lambda_j\mathcal M(\boldsymbol\lambda+\mathbf e_j,\mathbf k)
   +k_j\mathcal M(\boldsymbol\lambda+\mathbf e_j,\mathbf k-\mathbf e_j).
\]
The usual finite integration-by-parts formula in powers of $\Log R_{z_j}$
therefore supplies a normalized primitive in each $z_j$. For example,
with the other factors fixed and $\lambda_j\ne1$, a primitive at
$k_j=k$ is
\[
 \sum_{h=0}^k\frac{(-1)^hk!}{(k-h)!(\lambda_j-1)^{h+1}}
      \mathcal M(\boldsymbol\lambda-\mathbf e_j,
                         \mathbf k-h\mathbf e_j).
\]
When $\lambda_j=1$ it is the same moment with exponent $0$, logarithmic
degree $k+1$, divided by $k+1$. Subtract its value at a chosen base point
in the same half-plane to fix the additive constant. The proof is
coefficientwise differentiation of the completed integrals; no
uncompleted singular restriction is used.
